\slajdid{03-s004}
\begin{frame}[label=03-s004]{Izbor kola i neiskorišćeni ulazi}
% element: 03-s004:e01
U jednom SSI pakovanju po pravilu se nalaze \textbf{kola iste vrste}.
Za dva dvoulazna I kola, jedno troulazno I kolo i invertor početno su potrebna
\textbf{tri čipa}: sa I2 kolima, sa I3 kolima i sa invertorima.
\par\vspace{1mm}
% element: 03-s004:e02
\Rezultat{\textcolor{DEisticanje}{\textbf{Neiskorišćeni ulazi i ulazi neiskorišćenih delova čipa moraju biti na odgovarajućim neaktivnim nivoima.}}
U suprotnom mogu poremetiti rad aktivnog dela mreže.}
\par\vspace{1mm}
% element: 03-s004:e03
\textbf{Dva dvoulazna I kola mogu zameniti jedno troulazno:}
\[\textcolor{DEisticanje}{F=A\cdot B\cdot C=A\cdot(B\cdot C).}\]
Tako možemo koristiti manje čipova i manje vrsta čipova, ali samo ako je
\textbf{veće kašnjenje prihvatljivo}: signal prolazi kroz dva I2 kola umesto jednog I3 kola.
\par\vspace{1mm}
% element: 03-s004:e04
NI kolo sa spojenim ulazima može zameniti invertor: \(\overline{X\cdot X}=\overline X\).
\end{frame}
